% Oblique 3D sketch of the rotary (Furuta) pendulum: z up, arm in the xy plane.
\begin{tikzpicture}[x={(1cm,0cm)},y={(0.55cm,0.38cm)},z={(0cm,1cm)},>=Latex,scale=1.05,
  every node/.style={font=\small}]
\pgfmathsetmacro{\th}{-28}      % arm angle [deg]
\pgfmathsetmacro{\al}{28}       % pendulum angle [deg]
\pgfmathsetmacro{\Lr}{2.6}
\pgfmathsetmacro{\Lp}{2.9}
\pgfmathsetmacro{\tx}{\Lr*cos(\th)}
\pgfmathsetmacro{\ty}{\Lr*sin(\th)}
\pgfmathsetmacro{\ex}{\tx-\Lp*sin(\al)*sin(\th)}
\pgfmathsetmacro{\ey}{\ty+\Lp*sin(\al)*cos(\th)}
\pgfmathsetmacro{\ez}{\Lp*cos(\al)}
% base plate
\begin{scope}[canvas is xy plane at z=-2.9]
  \fill[gray!12] (-2.2,-2.6) rectangle (3.6,2.6);
  \draw[gray!50] (-2.2,-2.6) rectangle (3.6,2.6);
\end{scope}
% motor body (cylinder)
\begin{scope}[canvas is xy plane at z=-2.9] \path coordinate (b1) at (-0.55,0) coordinate (b2) at (0.55,0); \end{scope}
\fill[c1!25!gray!60] (-0.55,0,-2.9) -- (-0.55,0,-0.35) -- (0.55,0,-0.35) -- (0.55,0,-2.9) -- cycle;
\begin{scope}[canvas is xy plane at z=-2.9]
  \fill[c1!25!gray!60] (0,0) circle (0.55);
\end{scope}
\draw[gray!60!black] (-0.55,0,-2.9) -- (-0.55,0,-0.35) (0.55,0,-2.9) -- (0.55,0,-0.35);
\begin{scope}[canvas is xy plane at z=-0.35]
  \fill[c1!15!gray!40,draw=gray!60!black] (0,0) circle (0.55);
\end{scope}
\node[white,font=\footnotesize\bfseries] at (0,0,-1.6) {DC};
% hub
\fill[gray!45] (-0.18,0,-0.35) rectangle (0.18,0,0);
% swept circle of the arm tip and theta reference
\begin{scope}[canvas is xy plane at z=0]
  \draw[gray,densely dotted] (0,0) circle (\Lr);
  \draw[gray] (0,0) -- (\Lr+0.5,0) node[pos=1,right,black]{$\theta=0$};
  \draw[->,c4,thick] (2.0,0) arc (0:\th:2.0);
  \node[c4] at ({2.45*cos(\th/2)},{2.45*sin(\th/2)}) {$\theta$};
\end{scope}
% shadow of the pendulum on the plate
\draw[gray!35,line width=3.5pt,line cap=round] (0,0,-2.9) -- (\tx,\ty,-2.9);
\draw[gray!45,line width=2.2pt,line cap=round] (\tx,\ty,-2.9) -- (\ex,\ey,-2.9);
% upright reference at the arm tip
\draw[gray,dashed] (\tx,\ty,0) -- (\tx,\ty,\Lp+0.3);
% arm
\draw[line width=5pt,c1,line cap=round] (0,0,0) -- (\tx,\ty,0);
\node[c1,anchor=south east] at (-0.15,0,0.2) {brazo $L_r$};
% pendulum
\draw[line width=3.5pt,c2,line cap=round] (\tx,\ty,0) -- (\ex,\ey,\ez);
\fill[c2] (\ex,\ey,\ez) circle (0.12);
\fill[white,draw=black] (\tx,\ty,0) circle (0.07);
\node[c2,right] at ({0.5*(\tx+\ex)+0.15},{0.5*(\ty+\ey)},{0.5*\ez}) {péndulo $L_p$};
% alpha arc (in the plane of the pendulum)
\draw[->,c4,thick] plot[domain=0:\al,samples=12]
  ({\tx-1.3*sin(\x)*sin(\th)},{\ty+1.3*sin(\x)*cos(\th)},{1.3*cos(\x)});
\node[c4] at ({\tx-1.75*sin(\al/2)*sin(\th)},{\ty+1.75*sin(\al/2)*cos(\th)},{1.75*cos(\al/2)}) {$\alpha$};
% motor voltage
\draw[->,thick] (-2.0,0,-1.6) -- (-0.65,0,-1.6) node[pos=0,left]{$V$};
% axes
\draw[->,gray] (-1.9,-2.1,-2.9) -- (-0.9,-2.1,-2.9) node[right]{$x$};
\draw[->,gray] (-1.9,-2.1,-2.9) -- (-1.9,-1.1,-2.9) node[above]{$y$};
\draw[->,gray] (-1.9,-2.1,-2.9) -- (-1.9,-2.1,-1.9) node[above]{$z$};
\end{tikzpicture}
